\documentclass[12pt, a4paper]{article}

\usepackage[utf8]{inputenc}
\usepackage[T1]{fontenc}
\usepackage[italian]{babel}
\usepackage{geometry}
\geometry{
    a4paper,
    total={170mm,257mm},
    left=20mm,
    top=20mm,
}
\usepackage{array}
\usepackage{longtable} % Necessario per le tabelle multipagina

\begin{document}

\subsection*{Project Scoping Meeting Agenda}

\noindent
\textbf{Dipendente}: \\
\textbf{Data}: 1 maggio \\
\textbf{Orario}: 10:00 \\
\textbf{Luogo}: Meeting virtuale (Microsoft Teams)

\subsubsection*{Oggetto}
Primo incontro conoscitivo e di scoping con NexusLog S.p.A. per il progetto ``Shipping-on-the-Air''.

\subsubsection*{Partecipanti}
% [Sezione lasciata intenzionalmente vuota]
\vspace{1.5cm}

\subsubsection*{Svolgimento}
\renewcommand{\arraystretch}{1.5}
\begin{longtable}{|>{\raggedright\arraybackslash}p{0.25\textwidth}|c|>{\raggedright\arraybackslash}p{0.55\textwidth}|}
\hline
\textbf{Argomento in agenda} & \textbf{Orario} & \textbf{Note} \\
\hline
\endfirsthead

\hline
\textbf{Argomento in agenda} & \textbf{Orario}  & \textbf{Note} \\
\hline
\endhead

\hline
\endfoot

\hline
\endlastfoot

Introduzione & 10:00 - 10:30 & NexusLog è un operatore di primissimo piano a livello nazionale. Ci comunicano sin da subito il loro interesse verso l'approccio agile (Scrumban) della nostra startup, ritenuto ideale per gestire l'elevata incertezza tecnologica del progetto. \\
\hline
Scopo del meeting & 10:30 - 10:45 & Definire il perimetro del progetto ``Shipping-on-the-Air'', chiarire le ambiguità iniziali e valutare la fattibilità tecnica della soluzione software per la gestione dei droni. \\
\hline
Descrizione dello stato corrente e del problema & 10:45 - 11:30 & L'attuale modello logistico di NexusLog, basato esclusivamente sul trasporto terrestre, risulta ormai saturo. Il traffico urbano causa inefficienze sistematiche e ritardi imprevedibili, rendendo i costi operativi per l'ultimo miglio sproporzionati, soprattutto per la movimentazione di piccoli pacchi. \vspace{0.5em} \newline La necessità primaria dell'azienda è quindi quella di ampliare l'offerta dei propri servizi e, parallelamente, abbattere in modo drastico i costi operativi legati alle consegne di pacchi leggeri. Raggiungere questo duplice obiettivo permetterebbe a NexusLog di diventare altamente competitiva nel settore delle spedizioni ultra-rapide e di aggredire il mercato del \textit{food delivery} con soluzioni proprietarie. \\
\hline
Descrizione dello stato finale & 11:30 - 12:15 & Il committente vuole ampliare le consegne introducendo flotte di droni in città, con l'obiettivo di garantire spedizioni in meno di 2 ore per i pacchi e in meno di 1 ora per le consegne a domicilio di take away. \vspace{0.5em} \newline È stato però chiarito un punto fondamentale: pur volendo puntare forte sull'innovazione, NexusLog ci tiene a tutelare la reputazione di affidabilità costruita negli anni. Per questo motivo, l'azienda continuerà a gestire internamente tutta la parte operativa e commerciale, delegando all'esterno esclusivamente lo sviluppo del software e della tecnologia. I droni se li gestiscono in-house. \\
\hline
Discussione dei requisiti & 12:15 - 13:00 & I macro requisiti architetturali individuati sono:
\begin{itemize}
    \item \textbf{Gestione Ordini:} Ricezione e preparazione degli ordini.
    \item \textbf{Logistica e Flotta:} Verifica vincoli pre-volo (batteria, \textit{payload}).
    \item \textbf{Navigazione:} Calcolo rotte e algoritmi di volo.
    \item \textbf{Tracciamento Live:} Telemetria satellitare in tempo reale.
    \item \textbf{Notifiche \& Pagamenti:} Sistema \textit{event-driven} e gateway integrato.
\end{itemize}
I requisiti sono definiti ad alto livello. Saranno necessari meeting dedicati per ogni sottosistema. \\
\hline
\textbf{Pausa Pranzo} & \textbf{13:00 - 14:00} & Pausa pranzo.\\
\hline
Budget e tempi & 14:00 - 14:45 & Per mitigare l'esposizione finanziaria, il committente ha stanziato un budget limitato. Il progetto partirà come un \textit{Proof of Concept} (PoC) / \textit{Minimum Viable Product} (MVP). \vspace{0.5em} \newline NexusLog necessita di una prima versione funzionante entro 4 mesi per iniziare i test di volo sul campo nelle aree pilota. \\
\hline
Domande varie & 14:45 - 15:30 & 
\begin{itemize}
    \item Il \textit{procurement} hardware (droni, stazioni) sarà interamente gestito da NexusLog.
    \item Verranno fornite a breve le API proprietarie per interfacciarsi con la telemetria di bordo.
    \item Sarà necessaria un'integrazione software con i database attuali di NexusLog.
\end{itemize} \\
\hline
Documentazione & 15:30 - 16:00 & Non essendo chiari i dettagli algoritmici, non è possibile stendere la documentazione definitiva. Le specifiche di dettaglio verranno prodotte nei successivi incontri. \\
\end{longtable}

\subsubsection*{Conditions of Satisfaction}
\begin{itemize}
    \item Sviluppo di un \textit{Proof of Concept} funzionante per testare l'efficacia del servizio ``Shipping-on-the-Air'' prima di investimenti su larga scala.
    \item Rispetto rigoroso del budget limitato, concentrando le risorse esclusivamente sullo sviluppo software (\textit{outsourcing} tecnologico).
    \item Garanzia di abbattimento dei tempi di consegna.
    \item Preservazione della \textit{brand reputation} di NexusLog mediante un software affidabile, in modo da non impattare negativamente il \textit{core business} logistico tradizionale.
    \item Il sistema dovrà integrare con successo i 6 moduli cardine: gestione ordini, controllo flotta e \textit{payload}, \textit{routing} satellitare, tracciamento live, notifiche e pagamenti.
\end{itemize}

\subsubsection*{Pianificazione del prossimo meeting}
\begin{itemize}
    \item \textbf{Obiettivo:} Analisi del primo sottosistema (Gestione Ordini e Controllo Vincoli di Flotta/Payload) e revisione delle API fornite dal produttore dei droni.
    \item \textbf{Data:} 05/05/2026
    \item \textbf{Luogo:} Meeting Virtuale (Microsoft Teams)
    \item \textbf{Partecipanti:} Mario Rossi (NexusLog), e il nostro team di sviluppo software.
\end{itemize}

\vspace{1.5cm}
\noindent
\textbf{Firma}: \hrulefill \\ \\
\textbf{Data}: 01/05/2026

\end{document}